% GENERATED FILE. Edit canonical lessons, then run npm run generate:books.
% canonical-source-hash: fnv1a64:752275593934f9b1
% canonical-lessons: GE-C337-irren, GE-C337-irrtum, GE-C337-zweifel, GE-C337-zweifeln, GE-C337-bezweifeln, GE-R337-review-330-337

\chapter{Sich irren}
\label{ch:ge-sich-irren}
\hlchaptermodality{337}

\begin{chapteropening}
\textbf{By the end of this chapter:} \emph{I can say that someone is mistaken, express doubt and question a claim.}

Complete the last lesson of chapter 337: I can say that someone is mistaken, express doubt and question a claim.
\end{chapteropening}

\section[irren]{irren — to be mistaken (mostly reflexive, sich irren); to wander}

\label{lesson:GE-C337-irren}



\begin{quote}
Before the new one: say the German for on the one hand, then the German for on the other hand.
\end{quote}

\subsection*{You'll want to know: irren}
\textbf{irren} — \textquotedblleft{}to be mistaken (mostly reflexive, sich irren); to wander\textquotedblright{}.

\begin{quote}
Ich habe mich geirrt: Das Konzert ist am Sonntag, nicht am Samstag.
\end{quote}

\begin{quote}
\emph{I was mistaken: the concert is on Sunday, not Saturday.}
\end{quote}

Mostly reflexive: \emph{sich irren}, \emph{ich habe mich geirrt}. \emph{Irren ist menschlich} — to err is human.

\begin{cousinweb}
\emph{irren} is inherited from Old High German \emph{irren} and Proto-Germanic \emph{irzijaną} \textquotedblleft{}to delude, to mislead\textquotedblright{} (EtymDB), a cousin of Latin \emph{errāre} \textquotedblleft{}to wander\textquotedblright{}: to be mistaken is to wander off the path.
\end{cousinweb}

\begin{grammarlens}[title={What you've built}]
One more word for telling, explaining and describing at B1.
\end{grammarlens}

\subsection*{Your turn}
\begin{itemize}
\raggedright
  \item \emph{Say it:} \emph{irren}
  \item \emph{Say it:} \emph{irren}, once more
  \item \emph{Say it:} a sentence of your own with \emph{irren}
  \item \emph{Recall:} say the German for on the one hand, then the German for on the other hand, then say \emph{irren} again
\end{itemize}

\begin{tcolorbox}[breakable,colback=teal!4,colframe=teal!35!black,title={Before you move on}]
What does irren mean? (\textquotedblleft{}To be mistaken (mostly reflexive, sich irren); to wander\textquotedblright{}.) Say it once more.
\end{tcolorbox}
\section[der Irrtum]{der Irrtum — the mistake, the error}

\label{lesson:GE-C337-irrtum}



\begin{quote}
Before the new one: say the German for on the other hand, then the German for to be mistaken.
\end{quote}

\subsection*{You'll want to know: der Irrtum}
\textbf{der Irrtum} — \textquotedblleft{}the mistake, the error\textquotedblright{}.

\begin{quote}
Das war ein Irrtum: Die Rechnung ist schon bezahlt.
\end{quote}

\begin{quote}
\emph{That was an error: the bill has already been paid.}
\end{quote}

Plural: \emph{die Irrtümer}. \emph{Im Irrtum sein} — to be mistaken.

\begin{cousinweb}
\emph{Irrtum} is inherited through Middle High German \emph{irretuom} from Old High German \emph{irrituom} (EtymDB), built on \emph{irren} \textquotedblleft{}to err\textquotedblright{} with the noun ending \emph{-tum}.
\end{cousinweb}

\begin{grammarlens}[title={What you've built}]
One more word for telling, explaining and describing at B1.
\end{grammarlens}

\subsection*{Your turn}
\begin{itemize}
\raggedright
  \item \emph{Say it:} \emph{der Irrtum}
  \item \emph{Say it:} \emph{der Irrtum}, once more
  \item \emph{Say it:} a sentence of your own with \emph{der Irrtum}
  \item \emph{Recall:} say the German for on the other hand, then the German for to be mistaken, then say \emph{der Irrtum} again
\end{itemize}

\begin{tcolorbox}[breakable,colback=teal!4,colframe=teal!35!black,title={Before you move on}]
What does der Irrtum mean? (\textquotedblleft{}The mistake, the error\textquotedblright{}.) Say it once more.
\end{tcolorbox}
\section[der Zweifel]{der Zweifel — the doubt}

\label{lesson:GE-C337-zweifel}



\begin{quote}
Before the new one: say the German for to be mistaken, then the German for the mistake.
\end{quote}

\subsection*{You'll want to know: der Zweifel}
\textbf{der Zweifel} — \textquotedblleft{}the doubt\textquotedblright{}.

\begin{quote}
Ich habe keinen Zweifel: Sie ist die beste Spielerin.
\end{quote}

\begin{quote}
\emph{I have no doubt: she is the best player.}
\end{quote}

Plural: \emph{die Zweifel}. \emph{Ohne Zweifel} — without doubt.

\begin{cousinweb}
\emph{Zweifel} is inherited from Middle High German \emph{zwīvel}, Old High German \emph{zwīfal} and Proto-Germanic \emph{twīflaz} (EtymDB), built on the word for \textquotedblleft{}two\textquotedblright{}: a mind split in two.
\end{cousinweb}

\begin{grammarlens}[title={What you've built}]
One more word for telling, explaining and describing at B1.
\end{grammarlens}

\subsection*{Your turn}
\begin{itemize}
\raggedright
  \item \emph{Say it:} \emph{der Zweifel}
  \item \emph{Say it:} \emph{der Zweifel}, once more
  \item \emph{Say it:} a sentence of your own with \emph{der Zweifel}
  \item \emph{Recall:} say the German for to be mistaken, then the German for the mistake, then say \emph{der Zweifel} again
\end{itemize}

\begin{tcolorbox}[breakable,colback=teal!4,colframe=teal!35!black,title={Before you move on}]
What does der Zweifel mean? (\textquotedblleft{}The doubt\textquotedblright{}.) Say it once more.
\end{tcolorbox}
\section[zweifeln]{zweifeln — to have doubts}

\label{lesson:GE-C337-zweifeln}



\begin{quote}
Before the new one: say the German for the mistake, then the German for the doubt.
\end{quote}

\subsection*{You'll want to know: zweifeln}
\textbf{zweifeln} — \textquotedblleft{}to have doubts\textquotedblright{}.

\begin{quote}
Ich zweifle daran, dass der Bus pünktlich kommt.
\end{quote}

\begin{quote}
\emph{I have doubts that the bus will come on time.}
\end{quote}

\emph{An etwas zweifeln} — to have doubts about something. \emph{Ich zweifle} drops the \emph{e} of \emph{zweifele}.

\begin{cousinweb}
\emph{zweifeln} is inherited from Middle High German \emph{zwīvelen} (EtymDB), the verb of \emph{Zweifel} \textquotedblleft{}doubt\textquotedblright{}, which is built on the word for \textquotedblleft{}two\textquotedblright{}.
\end{cousinweb}

\begin{grammarlens}[title={What you've built}]
One more word for telling, explaining and describing at B1.
\end{grammarlens}

\subsection*{Your turn}
\begin{itemize}
\raggedright
  \item \emph{Say it:} \emph{zweifeln}
  \item \emph{Say it:} \emph{zweifeln}, once more
  \item \emph{Say it:} a sentence of your own with \emph{zweifeln}
  \item \emph{Recall:} say the German for the mistake, then the German for the doubt, then say \emph{zweifeln} again
\end{itemize}

\begin{tcolorbox}[breakable,colback=teal!4,colframe=teal!35!black,title={Before you move on}]
What does zweifeln mean? (\textquotedblleft{}To have doubts\textquotedblright{}.) Say it once more.
\end{tcolorbox}
\section[bezweifeln]{bezweifeln — to doubt (something), to question}

\label{lesson:GE-C337-bezweifeln}



\begin{quote}
Before the new one: say the German for the doubt, then the German for to have doubts.
\end{quote}

\subsection*{You'll want to know: bezweifeln}
\textbf{bezweifeln} — \textquotedblleft{}to doubt (something), to question\textquotedblright{}.

\begin{quote}
Der Arzt bezweifelt, dass das Medikament hilft.
\end{quote}

\begin{quote}
\emph{The doctor doubts that the medicine helps.}
\end{quote}

Inseparable, with a direct object: \emph{Ich bezweifle das.} \emph{Zweifeln} takes \emph{an}; \emph{bezweifeln} takes the thing itself.

\begin{cousinweb}
\emph{bezweifeln} is built with \emph{be-} on \emph{zweifeln}, the verb of \emph{Zweifel} \textquotedblleft{}doubt\textquotedblright{}, which goes back to Proto-Germanic \emph{twīflaz} (EtymDB).
\end{cousinweb}

\begin{grammarlens}[title={What you've built}]
One more word for telling, explaining and describing at B1.
\end{grammarlens}

\subsection*{Your turn}
\begin{itemize}
\raggedright
  \item \emph{Say it:} \emph{bezweifeln}
  \item \emph{Say it:} \emph{bezweifeln}, once more
  \item \emph{Say it:} a sentence of your own with \emph{bezweifeln}
  \item \emph{Recall:} say the German for the doubt, then the German for to have doubts, then say \emph{bezweifeln} again
\end{itemize}

\begin{tcolorbox}[breakable,colback=teal!4,colframe=teal!35!black,title={Before you move on}]
What does bezweifeln mean? (\textquotedblleft{}To doubt (something), to question\textquotedblright{}.) Say it once more.
\end{tcolorbox}
\section[(dialogue)]{Gründe und Meinungen}

\label{lesson:GE-R337-review-330-337}



\begin{quote}
Say the German for the standpoint and the view.
\end{quote}

\subsection*{Your turn}
Say these aloud:

\begin{itemize}
\raggedright
  \item the standpoint, the view, the understanding, the aspect, to advocate
  \item to convince, convincing, the conviction, to talk, to argue
  \item to give reasons for, the justification, you see, especially since, to explain in detail
  \item to contradict, the contradiction, the objection, to object, to deny
  \item to agree, the approval, agreed, to admit, to concede
  \item the discussion, the debate, the suggestion, to suggest, to stress
  \item the agreement, to give in, to weigh up, on the one hand, on the other hand
  \item to be mistaken, the mistake, the doubt, to have doubts, to doubt
\end{itemize}

\begin{tcolorbox}[breakable,colback=teal!4,colframe=teal!35!black,title={Before you move on}]
The standpoint? (\textbf{der Standpunkt}.) The view? (\textbf{die Ansicht}.) The understanding? (\textbf{die Auffassung}.) The aspect? (\textbf{der Gesichtspunkt}.) To advocate? (\textbf{vertreten}.) To convince? (\textbf{überzeugen}.) Convincing? (\textbf{überzeugend}.) The conviction? (\textbf{die Überzeugung}.) To talk? (\textbf{überreden}.) To argue? (\textbf{argumentieren}.) To give reasons for? (\textbf{begründen}.) The justification? (\textbf{die Begründung}.) You see? (\textbf{nämlich}.) Especially since? (\textbf{zumal}.) To explain in detail? (\textbf{erläutern}.) To contradict? (\textbf{widersprechen}.) The contradiction? (\textbf{der Widerspruch}.) The objection? (\textbf{der Einwand}.) To object? (\textbf{einwenden}.) To deny? (\textbf{bestreiten}.) To agree? (\textbf{zustimmen}.) The approval? (\textbf{die Zustimmung}.) Agreed? (\textbf{einverstanden}.) To admit? (\textbf{zugeben}.) To concede? (\textbf{einräumen}.) The discussion? (\textbf{die Diskussion}.) The debate? (\textbf{die Debatte}.) The suggestion? (\textbf{der Vorschlag}.) To suggest? (\textbf{vorschlagen}.) To stress? (\textbf{betonen}.) The agreement? (\textbf{die Einigung}.) To give in? (\textbf{nachgeben}.) To weigh up? (\textbf{abwägen}.) On the one hand? (\textbf{einerseits}.) On the other hand? (\textbf{andererseits}.) To be mistaken? (\textbf{irren}.) The mistake? (\textbf{der Irrtum}.) The doubt? (\textbf{der Zweifel}.) To have doubts? (\textbf{zweifeln}.) To doubt? (\textbf{bezweifeln}.)
\end{tcolorbox}
